\section{Dérivation bidirectionnelle de l'exposant électronique à partir du parcours enrichi}
\label{sec:cierre-electron-two-way-rev11}
\index{électron!lecteur bidirectionnel}
\index{calendrier d'ordre 108!exposant électronique}

Le parcours enrichi du point central donne deux représentations
complémentaires de sa mémoire : le mot tritique complet
\(\gamma_{108}\in\mathbb F_3^{108}\) et l'enregistrement dodécaphasique signé
\(\tau_{12}\in\{-1,0,+1\}^{12}\). Les deux lecteurs définis ci-dessous
sont calculés exclusivement à partir de ces représentations, du calendrier et des
constantes HMT déjà construites.

\subsection{Lecteur des déplacements sur le calendrier cyclique d'ordre 108}

Soit \(S\) le déplacement cyclique dans \(\mathbb Z/108\mathbb Z\). Pour
\(k\in\mathbb Z/108\mathbb Z\), nous définissons
\begin{equation}
 D_k(\Gamma):=
 \#\{t\in\mathbb Z/108\mathbb Z:\gamma_t\ne\gamma_{t+k}\}
 =\lVert(I-S^k)\gamma\rVert_0.
 \label{eq:lector-dk-ct108-rev11}
\end{equation}
Le triplet entier du lecteur est
\begin{equation}
 K_D(\Gamma)=
 \left(D_{27}+D_{36},\ D_1,\ \frac{D_4-D_3}{2}\right)
 =:(A_D,B_D,C_D),
 \label{eq:KD-rev11}
\end{equation}
et son évaluation scalaire est
\begin{equation}
 \kappa_D(\Gamma)=A_D-\alpha_{\rm HMT}B_D+\Delta_4 C_D.
 \label{eq:kappaD-rev11}
\end{equation}
Les déplacements (27) et (36) sont les relèvements chronologiques de
\(90^\circ\) et de \(120^\circ\) parce que neuf pas discrets représentent
\(30^\circ\). Le terme \(D_1\) mesure la frontière locale immédiate et la
différence \(D_4-D_3\) compare les deux généalogies avant leur quotient de
hauteur. Dans le domaine des 324 parcours, on a
\(\gamma_{t+54}=\gamma_t\) ; tous les \(D_k\) sont pairs et \(C_D\) est entier.

\subsection{Lecteur dodécaphasique des différences de corrélation}

Pour le mot \(\tau=(\tau_r)_{r\in\mathbb Z/12\mathbb Z}\), soient
\[
 C_k^{\rm corr}(\tau)=\sum_{r=0}^{11}\tau_r\tau_{r+k},
 \qquad
 A_\ell(\tau)=\sum_{r=0}^{11}
 \left(\sum_{j=0}^{\ell-1}\tau_{r+j}\right)^2.
\]
L'obstruction généalogique primitive et la mémoire antipodale sont
\begin{align}
 \Omega_{90\mid120}(\tau)
 &=4A_3-3A_4
 =-2C_1^{\rm corr}-4C_2^{\rm corr}-6C_3^{\rm corr},
 \label{eq:omega-two-way-rev11}\\
 P_6(\tau)&=\sum_{r=0}^{5}\tau_r\tau_{r+6}.
 \label{eq:p6-two-way-rev11}
\end{align}
Le symbole \(P_6\) évite la collision avec le cycle court du corridor des
constantes. Le second triplet et son évaluation sont
\begin{equation}
 K_\Omega(\Gamma)=\bigl(\Omega_{90\mid120}(\tau_\Gamma),54,
 P_6(\tau_\Gamma)\bigr),
 \qquad
 \kappa_\Omega=\Omega_{90\mid120}-54\alpha_{\rm HMT}+P_6\Delta_4.
 \label{eq:kappa-omega-rev11}
\end{equation}

\HMTClaim{MD-R-ELECTRON-OMEGA-PRIMITIVE}
\begin{lemma}[Unicité linéaire entière primitive]
\label{lem:unicidad-omega-rev11}
Soit \(L=uA_3+vA_4\), avec \(u,v\in\mathbb Z\), un contraste qui s'annule
lorsque les généalogies descendent à une même hauteur,
\(A_3=3a\), \(A_4=4a\). Alors
\((u,v)=n(4,-3)\). Le signe \(90-120\) étant fixé,
\(\Omega_{90\mid120}=4A_3-3A_4\) est le générateur primitif.
\end{lemma}

\begin{proof}
L'annulation pour tout \(a\) exige \(3u+4v=0\). Puisque
\(\gcd(3,4)=1\), les solutions entières sont
\((u,v)=n(4,-3)\). La primitivité correspond à \(|n|=1\).
\end{proof}

Le lemme sélectionne \(\Omega\) dans la classe des contrastes linéaires
entiers des énergies \(A_3,A_4\). Son domaine d'unicité est cette classe
fonctionnelle ; les invariants non linéaires relèvent d'un problème de sélection
distinct.

\subsection{Évaluation du parcours central}

Le parcours électronique \(\Gamma_e=(5,5,++)\) fournit, avant la
confrontation métrologique,
\[
 (D_1,D_3,D_4,D_{27},D_{36})=(54,56,68,48,32),
\]
et son enregistrement dodécaphasique signé est
\[
 \tau_e=(+,+,+,-,-,-,+,+,+,-,-,-).
\]
Par conséquent,
\[
 K_D(\Gamma_e)=(48+32,54,(68-56)/2)=(80,54,6),
\]
\[
 K_\Omega(\Gamma_e)=(80,54,6).
\]

\HMTClaim{MD-R-ELECTRON-TWOWAY-REDUCTION}
\begin{theorem}[Réduction électronique commune des lecteurs bidirectionnels]
\label{thm:reduccion-electron-two-way-rev11}
Les lecteurs \(K_D\) et \(K_\Omega\) coïncident sur le parcours central et
déterminent
\begin{equation}
 \kappa_e=80-54\alpha_{\rm HMT}+6\Delta_4.
 \label{eq:kappa-electron-demostrado-rev11}
\end{equation}
\end{theorem}

\begin{proof}
Les deux évaluations précédentes donnent le même triplet entier. La substitution
de ses composantes dans
\(A-\alpha_{\rm HMT}B+\Delta_4C\) prouve
\eqref{eq:kappa-electron-demostrado-rev11}. Toute quantité employée dans cette
réduction appartient au parcours, à ses deux enregistrements ou aux constantes HMT
préalablement construites.
\end{proof}

\paragraph{Réalisation ultérieure dans la carte énergétique.}
Avec \(R_{\rm act}=H_5/\etaRet\), l'action positive sur la valeur de base
\(\mathfrak e_0\) produit
\begin{equation}
 \mathfrak e_e^\beta
 =\mathfrak e_0R_{\rm act}^{\kappa_e}.
 \label{eq:masa-electron-adimensional-two-way-rev11}
\end{equation}
Après application de la carte actuelle vers les MeV, on obtient
\begin{equation}
 \varsigma_{u_E}(\mathfrak e_e^\beta)\big|_{u_E=1\,\mathrm{MeV}}
 =0.5109989506900008086082571648\ldots\ \mathrm{MeV}.
 \label{eq:masa-electron-two-way-rev11}
\end{equation}
La première égalité appartient au théorème interne conditionné par la valeur de base et
la droite d'action ; la seconde est sa réalisation dimensionnelle. La valeur
externe est réservée à la confrontation métrologique.

L'entier (80) procède de \(D_{27}+D_{36}\) et, dans le lecteur dodécaphasique,
de \(\Omega_{90\mid120}\). La coïncidence avec d'autres cardinaux internes est
un contrôle de compatibilité et ne remplace pas cette provenance opératorielle. Le
parcours d'orientation \((--)\) est une translation temporelle du parcours \((++)\)
et conserve le même profil, conformément à la parité de masse sous conjugaison.

\subsection{Séparation des lecteurs et des échelles de retour}

La coïncidence numérique des composantes internes reste subordonnée à leur
type. Au point central interviennent notamment les objets
\[
 R_{\rm atlas}(5,5)=(24,0,12,0,-12),
 \qquad
 L_{\kappa_0}(\Gamma_e^{++})=(24,0,0,0,-12),
\]
\[
 \begin{gathered}
 k_{\Delta_4}^{\rm dod},k_{\Delta_4}^{\rm pro}:
 \mathscr R_{324}\longrightarrow\mathbb Z,\\
 k_{\Delta_4}^{\rm dod}(\Gamma_e^{++})=12,
 \qquad
 k_{\Delta_4}^{\rm pro}(\Gamma_e^{++})=4,\\
 K_D(\Gamma_e^{++})=K_\Omega(\Gamma_e^{++})=(80,54,6).
 \end{gathered}
\]
Ici, \(\mathscr R_{324}\) désigne le domaine des 324 parcours orientés.
Le premier est une signature statique de famille ; le deuxième, un extracteur de
généalogie ; les troisième et quatrième termes sont des lecteurs torsionnels définis,
respectivement, par la lecture dodécaphasique et par le prolongement prospectif ;
le dernier objet est la paire de lecteurs bidirectionnels. Aucune égalité
partielle n'identifie ces morphismes ni ne permet de transférer un coefficient entre
lecteurs.

On distingue également quatre échelles temporelles : \(27\) est le premier retour
de classe, d'orbite et de cellule ; \(54=D_1(\Gamma_e)\) est le retour cinématique
observable dans le cycle ; \(108\) clôt le calendrier ; et \(216\) clôt l'état cinématique
orienté dans l'extension. Ce typage maintient distincts retour
positionnel, observation et holonomie orientée.

\paragraph{Chaîne causale du résultat électronique}
Le parcours et ses deux enregistrements déterminent \(K_D\), \(K_\Omega\), leur réduction
commune et l'évaluation de \(\kappa_e\). La réalisation en MeV applique ensuite
la carte énergétique ; la masse externe n'intervient que dans la confrontation. La paire
\(\mathbf K=(K_D,K_\Omega)\) est le résultat opératoriel complet de HMT sur
chaque parcours. Dans une fibre non centrale, sa représentation par un unique scalaire
physique exige un morphisme supplémentaire entre l'opérateur de fibre et la carte
dimensionnelle ; cette exigence ne modifie pas la loi interne et n'autorise pas à choisir
une cellule par proximité métrologique.
